\begingroup
\section{Foundation F57: stationary side mass and first moment}
\label{module:F57}
\noindent\textit{Audited source: \nolinkurl{convex_largeN_polya_szego_stage57_stationary_flux_reduction.tex}.}
\subsection{Active stationary polygons}

Let
\begin{equation}
 P=\bigcap_{i=1}^M\{x\in\mathbb R^2:x\cdot\nu_i<h_i\}
\end{equation}
be a bounded convex active polygon.  Write
\begin{equation}
 \nu_i=(\cos\theta_i,\sin\theta_i),\qquad
 \tau_i=(-\sin\theta_i,\cos\theta_i),
\end{equation}
and let $S_i$ be the side on the line $x\cdot\nu_i=h_i$.
Every point of $S_i$ has the unique representation
\begin{equation}\label{F57:eq:side-coordinate}
 x=h_i\nu_i+s\tau_i,\qquad s_i^-<s<s_i^+.
\end{equation}
Thus its length is
$\ell_i=s_i^+-s_i^-$.  Let $u>0$ be the $L^2(P)$-normalized first
eigenfunction,
\begin{equation}\label{F57:eq:eigen}
 -\Delta u=\lambda u\quad\hbox{in }P,\qquad
 u=0\quad\hbox{on }\partial P,\qquad
 \int_Pu^2=1,
\end{equation}
and put
\begin{equation}
 p_i(s)=-\partial_\nu u(h_i\nu_i+s\tau_i)\ge0.
\end{equation}
Convex polygonal regularity gives $u\in H^2(P)$ and
$p_i\in L^2(S_i)$.  The distributed Lipschitz-domain derivative, or
equivalently approximation inside a fixed active chart, gives the Hadamard
formula
\begin{equation}\label{F57:eq:Hadamard}
 D\lambda(P)[V]=-
 \sum_i\int_{S_i}p_i^2(V\cdot\nu_i)\,ds
\end{equation}
for every vertex-compatible polygonal material field $V$.

We use the scale-invariant functional
\begin{equation}
 \Lambda(P)=\frac{|P|}{\pi}\lambda_1(P).
\end{equation}
Stationarity of $\Lambda$ is equivalent to stationarity of $\lambda_1$
under the constraint $|P|=A$.

\subsection{The two exact sidewise moment identities}

\begin{theorem}[Stationary flux equidistribution]\label{F57:thm:flux}
Suppose that $P$ is stationary for $\Lambda$ with respect to all active
vertex perturbations.  Then, on every side $S_i$,
\begin{align}
 \frac1{\ell_i}\int_{s_i^-}^{s_i^+}p_i(s)^2\,ds
 &=\frac{\lambda}{|P|},\label{F57:eq:mean}\\
 \int_{s_i^-}^{s_i^+}
 \left(s-\frac{s_i^-+s_i^+}{2}\right)p_i(s)^2\,ds
 &=0.\label{F57:eq:first-moment}
\end{align}
Both identities are independent of the smallest side and of all mesh
ratios.
\end{theorem}

\begin{proof}
Work first with the Lagrangian
\begin{equation}
 \mathcal F(P)=\lambda_1(P)+\mu |P|.
\end{equation}

Fix the normal fan and replace $h_i$ by $h_i+t z_i$.  On the open side
$S_i$ the induced normal velocity is $z_i$; on every other open side it is
zero.  The endpoint motions are tangential to the neighboring sides and do
not contribute to the boundary integrals.  Consequently
\begin{align}
 D\lambda(P)[z]&=-\sum_i z_i\int_{S_i}p_i^2\,ds,
 &D|P|[z]&=\sum_i z_i\ell_i.
\end{align}
The $z_i$ are independent in an active chart.  Stationarity gives
\begin{equation}\label{F57:eq:multiplier-row}
 \int_{S_i}p_i^2\,ds=\mu\ell_i
 \qquad(1\le i\le M).
\end{equation}

To identify the multiplier, use the dilation field $V(x)=x$.  Scaling gives
\begin{equation}
 D\lambda(P)[x]=-2\lambda,\qquad D|P|[x]=2|P|.
\end{equation}
Hence $-2\lambda+2\mu|P|=0$, or
\begin{equation}\label{F57:eq:mu}
 \mu=\lambda/|P|.
\end{equation}
Equations \eqref{F57:eq:multiplier-row}--\eqref{F57:eq:mu} prove
\eqref{F57:eq:mean}.

It remains to rotate one supporting line.  Replace
$\theta_i$ by $\theta_i+tq_i$ while keeping $h_i$ fixed.  Differentiating
\begin{equation}
 x_t\cdot\nu_i(t)=h_i
\end{equation}
at a point $x=h_i\nu_i+s\tau_i$ gives
\begin{equation}\label{F57:eq:rot-velocity}
 V\cdot\nu_i=-s q_i.
\end{equation}
The normal velocity is zero on every other open side.  Therefore
\begin{align}
 D\lambda(P)[q_i]
 &=q_i\int_{s_i^-}^{s_i^+}s p_i(s)^2\,ds,\notag\\
 D|P|[q_i]
 &=-q_i\int_{s_i^-}^{s_i^+}s\,ds.
\end{align}
Stationarity of $\lambda+\mu|P|$ and the sign convention in
\eqref{F57:eq:multiplier-row} give
\begin{equation}\label{F57:eq:raw-first}
 \int_{s_i^-}^{s_i^+}s p_i(s)^2\,ds
 =\mu\int_{s_i^-}^{s_i^+}s\,ds.
\end{equation}
Subtract the midpoint times \eqref{F57:eq:multiplier-row} from
\eqref{F57:eq:raw-first}.  This is \eqref{F57:eq:first-moment}.
\end{proof}

\begin{corollary}[Sidewise probability law]\label{F57:cor:probability}
For every side define
\begin{equation}\label{F57:eq:probability}
 d\mathbb P_i(s)=
 \frac{p_i(s)^2}{(\lambda/|P|)\ell_i}\,ds.
\end{equation}
Then $\mathbb P_i$ is a probability measure and
\begin{equation}
 \int s\,d\mathbb P_i(s)=\frac{s_i^-+s_i^+}{2}.
\end{equation}
In particular, a stationary short side cannot be analyzed using only its
total flux: its first flux moment is fixed as well.
\end{corollary}

\begin{proof}
This is exactly \eqref{F57:eq:mean} and \eqref{F57:eq:first-moment}.
\end{proof}

\begin{remark}[Why the moment row is not redundant]
For a shallow corner model
$p(s)^2\simeq C s^{2\varepsilon_-}
(\ell-s)^{2\varepsilon_+}$, the normalized first moment is the mean of a
beta distribution.  Its displacement from $\ell/2$ has the sign of
$\varepsilon_--\varepsilon_+$.  Thus \eqref{F57:eq:first-moment} detects an
imbalance of the two endpoint corner exponents which the average identity
\eqref{F57:eq:mean} cannot see.
\end{remark}
\endgroup
